\metatitle{Polynomials indexed by Dyck paths}

\metadescription{One Dyck path determines a Ferrers board, a 312-avoiding
permutation and its Bruhat interval, a unit interval graph, a skew shape and a
lattice path matroid. We compare the polynomials attached to these objects:
rook-Eulerian and peak polynomials, acyclic orientations, natural descents,
non-nesting rooks and matroid Z-polynomials.}
\metakeywords{Dyck paths,area sequences,Ferrers boards,Bruhat order,312-avoiding
permutations,unit interval graphs,lattice path matroids,rook-Eulerian
polynomials,peaks,acyclic orientations,natural descents,Z-polynomials}


\section[dyckPathPolynomials]{Polynomials indexed by Dyck paths}

Many real-rooted families in algebraic combinatorics are indexed by a Catalan
object: a Dyck path, a Ferrers board inside a staircase, a $312$-avoiding
permutation, a unit interval graph or a skew shape.  This page fixes one
correspondence between these objects, so that the polynomials can be compared
directly.  Polynomials that look unrelated in one model are sometimes simple
statistics in another.


\subsection[dyckPathObjects]{One Dyck path, several objects}

Let $D$ be a \hyperref[qtCatalan]{Dyck path} of semilength $n$, with steps
$U$ and $D$.  Its \hyperref[areaSequence]{area sequence}
$a=(a_1,\dotsc,a_n)$ records the height of the path just before its $j$th
up-step, so $a_1=0$ and $a_{j+1}\leq a_j+1$.

\begin{itemize}
\item The \hyperref[unitIntervalGraph]{unit interval graph} $G_D$ on $[n]$
  joins each $j$ to its $a_j$ immediate predecessors $j-a_j,\dotsc,j-1$.
\item The \defin{Ferrers board}\label{dyckPathFerrersBoard} $\lambda_D$ has
  $\lambda_i$ boxes in row $i$, where $\lambda_i$ is the largest $j\geq i$ that
  is $i$ or a neighbor of $i$ in $G_D$.  The rows grow weakly from top to
  bottom, $\lambda_i\geq i$ and $\lambda_n=n$.  The boxes $(i,c)$ with $i<c$
  are the edges of $G_D$, and reading the upper right boundary of the board
  from the top left, with east steps as $U$ and south steps as $D$, gives the
  path $D$ back.
\item The permutation $\pi_D$ places a rook in each row of $\lambda_D$, from
  the top, as far to the right as possible without attacking the rooks above.
  It avoids the pattern $312$, every $312$-avoiding permutation arises once,
  and the permutations $\sigma$ with $\sigma_i\leq\lambda_i$ for all $i$, that
  is, the complete rook placements on $\lambda_D$, form the lower
  \hyperref[bruhatOrder]{Bruhat interval} $[\mathrm{id},\pi_D]$
  \cite[Prop.~3]{AlexanderssonJalQuemener2025x}.
\item The \defin{peak intervals}\label{dyckPathPeakIntervals} are the intervals
  $[i,\lambda_i]$ for the rows $i$ with $\lambda_i>\lambda_{i-1}$, where
  $\lambda_0=0$.  They correspond to the peaks of $D$, and they are the maximal
  cliques of $G_D$: if $i<j$ then $i$ and $j$ are adjacent exactly when
  $j\leq\lambda_i$, so a clique with smallest element $i$ lies in
  $[i,\lambda_i]$.
\item Number the peak intervals $[\ell_1,u_1],\dotsc,[\ell_r,u_r]$ from left to
  right.  The \hyperref[latticePathMatroids]{lattice path matroid} $M_D$ on
  $[n]$ has as bases the sets $\{b_1<\dotsb<b_r\}$ with
  $\ell_k\leq b_k\leq u_k$.  It is the lattice path matroid of the
  \hyperref[skewShapes]{skew shape} $F_D$ with $r$ rows, whose $k$th row from
  the bottom occupies the columns $\ell_k-k+1,\dotsc,u_k-k$.  The entries of a
  peak interval are the possible positions of the north step of that row in a
  lattice path through $F_D$.  A peak interval with a single element gives an
  empty row, and a coloop of $M_D$.
\end{itemize}

Conversely, let $\lambda/\mu$ be a skew shape with $r$ rows, empty rows
allowed, with $\lambda_1=c$, $\mu_r=0$ and $\mu_j\leq\lambda_{j+1}$ for all
$j$; these are the shapes whose lattice path matroid has no loops.  Each such
shape arises from exactly one Dyck path of semilength $r+c$.  Straight shapes
correspond to $\ell_k=k$ for all $k$.

\begin{example}
Let $a=(0,1,2,2,3)$, so $D=UUUDUUDDDD$.  The board is
$\lambda_D=(3,5,5,5,5)$ and $\pi_D=35421$.  The graph $G_D$ has the eight edges
$12$, $13$, $23$, $24$, $25$, $34$, $35$ and $45$, and its maximal cliques are
the peak intervals $A_1=[1,3]$ and $A_2=[2,5]$.  The skew shape is $F_D=(3,2)$:
its bottom row admits the north-step positions $1,2,3$ and its top row the
positions $2,3,4,5$.  The matroid $M_D$ has rank two, and its bases are all
pairs from $[5]$ except $\{4,5\}$.
\begin{figure}
\svgimg[width=0.95\textwidth]{./svg-images/dyck-path-objects.svg}{The board (3,5,5,5,5) with the rooks of the permutation 35421, the unit interval graph with its two maximal cliques, and the skew shape (3,2) with the possible north-step positions in each row.}
\end{figure}
\end{example}


\subsection[dyckPathBruhatPolynomials]{Statistics on the Bruhat interval}

The complete rook placements on $\lambda_D$, equivalently the permutations in
$[\mathrm{id},\pi_D]$, carry three natural statistics.  Write
$[k]_q=1+q+\dotsb+q^{k-1}$.

The number of inversions is the length in the Bruhat order, and
\[
  X_D(q)\coloneqq\sum_{\sigma\leq\pi_D}q^{\operatorname{inv}(\sigma)}
  =\prod_{j=1}^{n}[a_j+1]_q.
\]
This is the \hyperref[aoRooks]{rook model} for the weight
$X_D(q)$ on the page on \hyperref[unitIntervalAO]{Touchard--Riordan
polynomials}.  The same product counts the
\hyperref[aoOnUnitIntervalGraphs]{acyclic orientations} of $G_D$ by ascents
\cite[Sec.~9]{AlexanderssonPanova2018}: inserting vertex $j$ into an acyclic
orientation of the earlier vertices, its $a_j$ earlier neighbors form a clique
that is already ordered, and the $a_j+1$ slots for $j$ give $0,\dotsc,a_j$ new
ascents.  It also counts the perfect matchings of $[2n]$ whose openers are at
the up-steps of $D$, by crossings, so summing over all Dyck paths gives the
\hyperref[touchardRiordanPolynomialDefinition]{Touchard--Riordan polynomial}
$T_n(q)$.  Explicitly, read the closers of such a matching $M$ from left to
right and record, for each closer, the index of its opener among the openers.
This gives a bijection with $[\mathrm{id},\pi_D]$, under which
$\operatorname{cr}(M)=a_1+\dotsb+a_n-\operatorname{inv}(\sigma)$.

The \defin{rook-Eulerian polynomial}\label{dyckPathRookEulerian} counts the
same permutations by ascents,
\[
  Q_D(t)\coloneqq\sum_{\sigma\leq\pi_D}t^{\operatorname{asc}(\sigma)}.
\]
For the full board $\lambda_D=(n,\dotsc,n)$ this is the
\hyperref[eulerianPolynomial]{Eulerian polynomial}.

\begin{theorem}[Alexandersson--Jal--Quemener, {\cite[Thms.~12--13]{AlexanderssonJalQuemener2025x}}]
For every Dyck path $D$, the polynomial $Q_D(t)$ is real-rooted.  More
generally this holds for row-complete rook placements on any Ferrers board with
$\lambda_i\geq i$, and the multivariate ascent-set enumerator is same-phase
stable.
\end{theorem}

The \defin{peak polynomial}\label{dyckPathPeakPolynomial} counts them by
interior peaks, positions $1<i<n$ with $\sigma_{i-1}<\sigma_i>\sigma_{i+1}$,
\[
  P_D(t)\coloneqq\sum_{\sigma\leq\pi_D}t^{\operatorname{pk}(\sigma)}.
\]

\begin{theorem}[Alexandersson--Jal--Quemener, \cite{AlexanderssonJalQuemener2025x}]
For every Ferrers board with weakly increasing rows and $\lambda_i\geq i$, the
peak polynomial of its row-complete rook placements is real-rooted.  In
particular $P_D(t)$ is real-rooted.
\end{theorem}

The inversions can be kept as well: for every $q\geq0$,
\[
  \sum_{\sigma\leq\pi_D}q^{\operatorname{inv}(\sigma)}t^{\operatorname{pk}(\sigma)}
\]
is real-rooted in $t$.  Indeed, if $\sigma_1=k$, then $\sigma$ has $k-1$
inversions involving its first entry, and removing the first row and column $k$
leaves a placement on a smaller board of the same kind.  So the inversion
weight multiplies each refinement by the first entry, in the induction of
\name{Per Alexandersson}, \name{Aryaman Jal} and \name{Maena Quemener}, by a
positive scalar, which preserves its interlacing conditions.  By the matching
bijection above, the peaks of the closer word, weighted by $q$ to the number
of crossings, therefore also give a real-rooted polynomial for each set of
openers.

Signing the entries connects peaks and descents.  Give each entry of
$\sigma$ a sign, and count ordinary descents of the signed word in the usual
order of the integers.  Each comparison of adjacent entries is decided by the
sign of the entry with larger absolute value.  An interior peak owns two
comparisons and gives exactly one descent for either sign, an entry owning one
comparison gives $1+x$ over its two signs, and the remaining $p+1$ entries own
none.  Hence a word with $p$ peaks gives
$\sum_{\epsilon}x^{\operatorname{des}(\epsilon\sigma)}
=2^{2p+1}x^p(1+x)^{n-2p-1}$, and
\[
  \sum_{\sigma\leq\pi_D}\ \sum_{\epsilon\in\{\pm1\}^n}
  x^{\operatorname{des}(\epsilon\sigma)}
  =2(1+x)^{n-1}\,P_D\Bigl(\frac{4x}{(1+x)^2}\Bigr).
\]
The same argument is used for the
\hyperref[schubertMatroidZPermutations]{signed-descent formula} for
$Z$-polynomials of Schubert matroids.

Keeping the whole descent set instead of the number of descents gives the
quasisymmetric function $\sum_{\sigma\leq\pi_D}\gessel_{n,\DES(\sigma)}$.  It is
not symmetric unless $a=(0,\dotsc,0)$ or $a=(0,1,\dotsc,n-1)$, but it is
positive in the quasisymmetric power-sum basis, by a sign-reversing involution;
see \hyperref[bruhatIntervalPsiPositive]{lower Bruhat intervals} on the page on
$P$-partitions.


\subsection[dyckPathGraphPolynomials]{Chromatic quasisymmetric functions}

The graph $G_D$ is the natural unit interval graph $H_h$ with Hessenberg
function $h=\lambda_D$, so the Dyck path also indexes its
\hyperref[chromaticQuasisymmetric]{chromatic quasisymmetric function}
\[
  \chrom_{G_D}(\xvec;q)=\sum_{\lambda}c_\lambda(q)\,\elementaryE_\lambda(\xvec)
\]
of \name{John Shareshian} and \name{Michelle L. Wachs}
\cite{ShareshianWachs2016}.  It is symmetric and Schur-positive, it is
elementary-positive at $q=1$ by \name{Tatsuyuki Hikita}'s proof of the
\hyperref[stanleyStembridgeConjecture]{Stanley--Stembridge conjecture}
\cite{Hikita2024x}, and it is related to the cohomology of regular
semisimple \hyperref[chromaticQuasisymmetricHessenberg]{Hessenberg
varieties}.  Several of the polynomials on this page are specializations of
it.

A sink of an orientation is a vertex with no outgoing edge.  Let
\[
  S_D(t;q)\coloneqq\sum_{O}q^{\operatorname{asc}(O)}t^{\operatorname{sink}(O)},
  \qquad S_D(t)\coloneqq S_D(t;1),
\]
summed over the acyclic orientations $O$ of $G_D$.  By the Shareshian--Wachs
sink identity, $S_D(t;q)=\sum_\lambda c_\lambda(q)t^{\ell(\lambda)}$
\cite[Thm.~5.3]{ShareshianWachs2016}, and setting $t=1$ gives $X_D(q)$.  This
\hyperref[chromaticSinkRealRooted]{sink specialization} is real-rooted in $t$
for every $q\geq0$ \cite[Thm.~5.2]{AlexanderssonLeite2026x}.

A half-sink is a vertex with no outgoing edge to a larger label.  Inserting the
vertices in decreasing order gives the product
\[
  \sum_{O}q^{\operatorname{asc}(O)}t^{\operatorname{hsink}(O)}
  =\prod_{i=1}^{n}\bigl(t+q+q^2+\dotsb+q^{\lambda_i-i}\bigr),
\]
over acyclic orientations, read off from the rows of the board
\cite[Prop.~40]{AlexanderssonLeite2026x}.

The \hyperref[gesselLengthSpecialization]{fundamental-length
specialization} gives \name{Francesco Brenti}'s
\hyperref[chromaticWPolynomial]{$w$-polynomial}
\[
  w_D(t)=(1-t)^{n+1}\sum_{m\geq0}\chi_{G_D}(m)\,t^m,\qquad
  \chi_{G_D}(m)=\prod_{j=1}^{n}(m-a_j),
\]
where the chromatic polynomial factors since the earlier neighbors of each
vertex form a clique.  It is real-rooted since $G_D$ is chordal
\cite[Thms.~4.11 and~4.15]{Brenti1992}.

\begin{example}
For $a=(0,1,2,2,3)$, every acyclic orientation has a sink, and the only sets of
two sinks are $\{1,4\}$ and $\{1,5\}$.  Each leaves a triangle with six acyclic
orientations, so $S_D(t)=60t+12t^2$, with $X_D(1)=72$ orientations in total.
The half-sink polynomial at $q=1$ is $t(t+1)(t+2)^2(t+3)$, and
$w_D(t)=48t^4+72t^5$.
\end{example}


\subsection[dyckPathSkewPolynomials]{Natural descents and non-nesting rooks}

Two polynomials are attached to the skew shape $F_D$.  For a
\hyperref[standardYoungTableau]{standard Young tableau} $T$ of shape $F_D$, a
\defin{natural descent} is an entry $j$ such that $j+1$ lies in a strictly
higher row.  Let
\[
  N_D(t)\coloneqq\sum_{T}t^{\operatorname{ndes}(T)},\qquad
  R_D(t)\coloneqq\sum_{k}r_k\,t^k,
\]
where $r_k$ is the number of
\hyperref[nonNestingRookPlacement]{non-nesting rook placements} on $F_D$ with
$k$ rooks.

Both polynomials are \hyperref[pEulerianPolynomial]{$P$-Eulerian
polynomials}.  Reading the cells row by row labels the cell poset
of $F_D$ naturally, and $N_D$ is its $P$-Eulerian polynomial.  The
non-nesting rook polynomial is the $P$-Eulerian polynomial of the width-two
poset of the boundary steps, $R_D=W_P$ \cite{AlexanderssonJal2024x}.  The
order ideals of the cell poset are the lattice paths through $F_D$, that is,
the bases of $M_D$, and the maximal cells of an order ideal form a non-nesting
rook placement.  Hence $R_D(1)$ is the number of bases of $M_D$.

Two further descriptions of $N_D$ follow.  Since the cell poset has dimension
two, label the cells row by row and read the labels column by column, from the
top, to get a permutation $w_D$.  The tableaux of shape $F_D$, written as words
in these labels, are exactly the permutations below $w_D$ in the right
\hyperref[weakOrder]{weak order}, with natural descents as ordinary descents.
This is a weak-order interval, in contrast to the Bruhat interval behind
$Q_D$ and $P_D$.  Moreover, $N_D$ is the $h$-polynomial of the order complex
of the bases of $M_D$ in Gale order, with the smallest and largest basis
removed:
\[
  N_D(t)=(1-t)^{m-1}f\Bigl(\frac{t}{1-t}\Bigr),
\]
where $f(z)$ counts the chains of these bases by size, and $m=|F_D|$.

\begin{proposition}
If $F_D$ has $m$ cells, then
\[
  [t]N_D(t)=R_D(1)-m-1.
\]
\end{proposition}

\begin{proof}
By the theory of $P$-partitions, $N_D(t)/(1-t)^{m+1}=\sum_{k\geq0}c_kt^k$,
where $c_k$ counts the order-preserving maps from the cell poset to
$\{0,\dotsc,k\}$.  Such maps with $k=1$ are the order ideals, so
$c_1=R_D(1)$, and comparing the coefficients of $t$ gives
$[t]N_D=c_1-(m+1)$.
\end{proof}

For two-row shapes there is an identity of whole polynomials after changing
the board.  If the top row has $s$ cells, the bottom row has $b$ cells, and the
top row starts $d\geq0$ columns to the right of the bottom row, then the
natural-descent polynomial equals the non-nesting rook polynomial of the
Ferrers board with rows $\theta_j=\min(s,s+d-j)$, for $1\leq j\leq b$.
A bijection records a tableau by its word in the letters $A$ (top row)
and $B$ (bottom row).  At each occurrence of $BA$, let $x$ and $y$ count the
letters $A$ and $B$ up to and including that $B$, and place a rook in row $y$
and column $s-x$; between two such occurrences, the word is a run of $A$
followed by a run of $B$.  For example, the shape $(3,2)$ gives the board $(2,1)$, and both
polynomials are $1+3t+t^2$.

Real-rootedness of natural-descent polynomials of straight shapes is
\hyperref[conjSYTNaturalDescents]{conjectured}.  The non-nesting rook numbers
are ultra-log-concave, but $R_D(t)$ need not be real-rooted
\cite{AlexanderssonJal2024x}.


\subsubsection[dyckPathParticles]{Tableaux as particle hops}

The skew shape $F_D$ can be found inside the board.  The cells $(i,c)$ of
$\lambda_D$ with $i<c$ are the edges of $G_D$.  Keep only the rows
$\ell_1,\dotsc,\ell_r$ where a peak interval starts; row $\ell_k$ contains the
edges $\{\ell_k,p\}$ with $\ell_k<p\leq u_k$.  Shifting the $k$th kept row
$k$ steps to the left, and stacking the rows from the bottom, gives $F_D$.

This leads to a particle model.  Place $r$ particles on the sites
$\ell_1<\dotsb<\ell_r$ of $[n]$.  In each step one particle moves one site to
the right, onto an empty site, and particle $k$ stops at $u_k$, so it never
leaves the clique $A_k$.  The positions at any moment form a basis of $M_D$,
and a complete run is a maximal chain of bases in Gale order.  A standard
Young tableau $T$ of shape $F_D$ gives a run: if $t$ lies in the $k$th row from
the bottom, the $t$th step moves particle $k$, and the $c$th cell of that row is the hop
of particle $k$ onto site $\ell_k+c$, that is, the edge $\{\ell_k,\ell_k+c\}$.  The columns of $T$ increase precisely when
no two particles meet, and every run arises once.  A natural descent $j$ of
$T$ is a time where the particle moved at step $j+1$ lies to the right of the
particle moved at step $j$.  Hence $N_D(t)$ counts the runs by the ascents of the word
recording which particle moves, and $R_D(1)$ counts the configurations.

\begin{example}
For $a=(0,1,2,2,3)$, particle $1$ moves from site $1$ to site $3$ and particle
$2$ from site $2$ to site $5$.  The tableau with rows $124$ and $35$ moves
particle $2$, particle $2$, particle $1$, particle $2$ and particle $1$.  Its
only natural descent is $3$, where the next moving particle lies to the right.
\begin{figure}
\svgimg[width=0.9\textwidth]{./svg-images/dyck-path-particles.svg}{The peak rows of the area of the Dyck path with area sequence 0,1,2,2,3, the tableau of shape (3,2) with rows 124 and 35, and the corresponding particle positions after each hop.}
\end{figure}
\end{example}

For a straight shape $\lambda$, the particles start at $1,\dotsc,r$, and a run
is a ballot sequence recording the row of each entry.  The natural descents
are then the descents of the ballot sequence, with rows numbered from the top.
For the two-row shape $(m,m)$ these are the valleys of a Dyck path of
semilength $m$, which have the \hyperref[narayanaNumbers]{Narayana}
distribution.


\subsection[dyckPathZPolynomials]{Matroid Z-polynomials}

Let $Z_D(t)$ be the \hyperref[matroidZPolynomial]{$Z$-polynomial} of the
lattice path matroid $M_D$.  For connected $F_D$ it is given by the
\hyperref[schubertMatroidZRookMotzkin]{Motzkin path and rook formula} on the
$2\times2$ squares of $F_D$.  In general $Z_D$ is the product over the
connected components of $F_D$, with a factor $1+t$ for each coloop, and for straight shapes it is the $Z$-polynomial
of a \hyperref[schubertMatroidZ]{Schubert matroid}.

Dyck paths enter a second time, through the
\hyperref[schubertMatroidZWordFormula]{word formula}.  The words
$a_1\dotsm a_n$ there are exactly the area sequences $(0,a_1,\dotsc,a_n)$ of the
Dyck paths $E$ of semilength $n+1$.  Since $[k]_q$ vanishes at $q=-1$ precisely
when $k$ is even, the number of odd letters is the order of vanishing of $X_E$
at $q=-1$.  For a Schubert matroid of rank $n$ with board of squares $B$,
\[
  Z(t)=\sum_{E}t^{\operatorname{ord}_{q=-1}X_E(q)}\,
  W_B\bigl(\{j\in[n-1]: a_{j+1}(E)\geq2\}\bigr),
\]
summed over Dyck paths $E$ of semilength $n+1$, where $a_{j+1}(E)$ is the
height of $E$ before its $(j+2)$nd up-step.  When the board of squares is empty,
for instance for the uniform matroid $U_{n,n+1}$, every $W_B$ equals one, and
\[
  N_{n+1}(t)=\sum_{E}t^{\operatorname{ord}_{q=-1}X_E(q)}
\]
is the \hyperref[narayanaNumbers]{Narayana polynomial}.  For $n=2$, the five
area sequences $000$, $001$, $010$, $011$ and $012$ give $1+3t+t^2$.


\subsection[dyckPathSemilengthFour]{The fourteen Dyck paths of semilength four}

\begin{figure}
\svgimg[width=0.95\textwidth]{./svg-images/dyck-paths-four.svg}{The fourteen Dyck paths of semilength four, labelled by their area sequences.}
\end{figure}

The first table lists the objects attached to each Dyck path: the board
$\lambda_D$, the permutation $\pi_D$, the peak intervals and the skew shape
$F_D$, written as $\lambda/\mu$ with rows from top to bottom.

\begin{tabular}{lllll}
\toprule
$a$ & $\lambda_D$ & $\pi_D$ & $A_1,\dotsc,A_r$ & $F_D$ \\
\midrule
$0000$ & $1234$ & $1234$ & $[1],[2],[3],[4]$ & $(0,0,0,0)$ \\
$0001$ & $1244$ & $1243$ & $[1],[2],[3,4]$ & $(1,0,0)$ \\
$0010$ & $1334$ & $1324$ & $[1],[2,3],[4]$ & $(1,1,0)/(1,0,0)$ \\
$0011$ & $1344$ & $1342$ & $[1],[2,3],[3,4]$ & $(1,1,0)$ \\
$0012$ & $1444$ & $1432$ & $[1],[2,4]$ & $(2,0)$ \\
$0100$ & $2234$ & $2134$ & $[1,2],[3],[4]$ & $(1,1,1)/(1,1,0)$ \\
$0101$ & $2244$ & $2143$ & $[1,2],[3,4]$ & $(2,1)/(1,0)$ \\
$0110$ & $2334$ & $2314$ & $[1,2],[2,3],[4]$ & $(1,1,1)/(1,0,0)$ \\
$0111$ & $2344$ & $2341$ & $[1,2],[2,3],[3,4]$ & $(1,1,1)$ \\
$0112$ & $2444$ & $2431$ & $[1,2],[2,4]$ & $(2,1)$ \\
$0120$ & $3334$ & $3214$ & $[1,3],[4]$ & $(2,2)/(2,0)$ \\
$0121$ & $3344$ & $3241$ & $[1,3],[3,4]$ & $(2,2)/(1,0)$ \\
$0122$ & $3444$ & $3421$ & $[1,3],[2,4]$ & $(2,2)$ \\
$0123$ & $4444$ & $4321$ & $[1,4]$ & $(3)$ \\
\bottomrule
\end{tabular}

The second table lists the polynomials.  The product $X_D(q)$ can be read off
from $a$ directly.

\begin{tabular}{lllllll}
\toprule
$a$ & $Q_D$ & $P_D$ & $S_D$ & $N_D$ & $R_D$ & $Z_D$ \\
\midrule
$0000$ & $t^3$ & $1$ & $t^4$ & $1$ & $1$ & $1+4t+6t^2+4t^3+t^4$ \\
$0001$ & $t^2+t^3$ & $1+t$ & $2t^3$ & $1$ & $1+t$ & $1+3t+3t^2+t^3$ \\
$0010$ & $t^2+t^3$ & $1+t$ & $2t^3$ & $1$ & $1+t$ & $1+3t+3t^2+t^3$ \\
$0011$ & $3t^2+t^3$ & $1+3t$ & $3t^2+t^3$ & $1$ & $1+2t$ & $1+4t+4t^2+t^3$ \\
$0012$ & $t+4t^2+t^3$ & $1+5t$ & $6t^2$ & $1$ & $1+2t$ & $1+2t+t^2$ \\
$0100$ & $t^2+t^3$ & $2$ & $2t^3$ & $1$ & $1+t$ & $1+3t+3t^2+t^3$ \\
$0101$ & $t+2t^2+t^3$ & $2+2t$ & $4t^2$ & $1+t$ & $1+2t+t^2$ & $1+2t+t^2$ \\
$0110$ & $3t^2+t^3$ & $2+2t$ & $3t^2+t^3$ & $1$ & $1+2t$ & $1+4t+4t^2+t^3$ \\
$0111$ & $t+6t^2+t^3$ & $2+6t$ & $4t+4t^2$ & $1$ & $1+3t$ & $1+6t+6t^2+t^3$ \\
$0112$ & $3t+8t^2+t^3$ & $2+10t$ & $8t+4t^2$ & $1+t$ & $1+3t+t^2$ & $1+3t+t^2$ \\
$0120$ & $t+4t^2+t^3$ & $4+2t$ & $6t^2$ & $1$ & $1+2t$ & $1+2t+t^2$ \\
$0121$ & $4t+7t^2+t^3$ & $4+8t$ & $8t+4t^2$ & $1+t$ & $1+3t+t^2$ & $1+3t+t^2$ \\
$0122$ & $7t+10t^2+t^3$ & $4+14t$ & $16t+2t^2$ & $1+t$ & $1+4t+t^2$ & $1+4t+t^2$ \\
$0123$ & $1+11t+11t^2+t^3$ & $8+16t$ & $24t$ & $1$ & $1+3t$ & $1+t$ \\
\bottomrule
\end{tabular}

Some polynomials coincide.  Reversing the Dyck path, as for $0001$ and
$0100$, gives isomorphic graphs, rotated skew shapes and isomorphic matroids,
so $S_D$, $N_D$, $R_D$ and $Z_D$ are unchanged.  The Bruhat interval does not
have this symmetry: $Q_D$ often agrees but $P_D$ usually differs, as for
$0001$ and $0100$.  The rows $0101$, $0112$, $0121$ and $0122$, whose shapes
have two nonempty rows, have $R_D=Z_D$, while the row $0111$ shows that the two
differ in general.
